% Independently audited by root and maps. Integral certificate in companion .py.
\paragraph{Theorem B.6.20 (Proved).}
For the marked fixed-curve surgery of the minimal $(3,4)$ family in
B.6.16, let $Q$ be the six-dimensional free-circle quotient of the
parameter mapping torus. Its primitive Euler class $e$ satisfies
$|\langle e^3,[Q]\rangle|=1$. In the oriented convention of B.6.17,
$\ell=-4$, hence the spin congruence of B.6.15 gives $k=-1$.
Consequently the spin-trace calculation of B.6.19 gives the zero class
in $\Theta_7$ for the original minimal-family monodromy.

\emph{Proof.}
Let $g,u,w,k$ be the integral fibre lattice of the free-circle quotient,
with dual forms $G,U,W,K$.  In this lattice notation the last
letter is the parameter direction, distinct from the integer cup
coefficient $k$ in the statement.  The four homology matrices are
\[
 A_1=\begin{pmatrix}1&0&0&0\\6&0&1&-1\\-6&-1&-1&1\\0&0&0&1\end{pmatrix},\quad
 A_2=\begin{pmatrix}1&0&0&0\\0&0&-1&0\\-6&1&0&1\\0&0&0&1\end{pmatrix},
\]
\[
 C_1=I+w(u^*-g^*),\qquad C_2=I+wg^*,\qquad A_1A_2C_1C_2=I.
\]
The finite caps have orders $3,4$ and power vectors
$v_1=(1,2,-4,0)$, $v_2=(-1,-3,3,0)$.
The other two caps are the ordinary primitive nodal products after
quotienting by the free transverse circle. The exact normal-torus
Mayer--Vietoris calculation in \texttt{manifolds\_mapping\_quotient.py} gives
$H_2(Q)=\mathbb Z$, and the regular $gk$-torus maps to twice its generator.
The primitive degree-two class $x$ therefore has fibre restriction
$E_0=2G\wedge K$ after choosing the sign of $x$.

\paragraph{Closed forms on the exterior.}
Remove four discs from the oriented base sphere and choose three disjoint
arcs from the fourth boundary to the first three.  A flat torus bundle over
this surface retracts over the resulting graph. Fibrewise constant forms
and forms supported in disjoint narrow strips about these arcs give a
multiplicative de Rham model: the fibre part is the exterior algebra of
$G,U,W,K$; the first three meridian parts are one-forms with values in this
algebra, with differential $A_i^T-I$.  This assertion can also be checked
by cutting the base to a disc, using the actual linear torus gluings, and
applying de Rham Mayer--Vietoris. In particular, representatives whose
mixed base terms have disjoint support have no product containing two
such base terms.

The integral normal-torus calculation gives the following representative
of $x$ on the exterior:
\[
 E_0=2GK,\qquad
 \eta_1=-4G-W,\quad\eta_2=-3G+U+K,\quad\eta_3=0.
\]
Here a boundary path with parameter $t$ carries the two-form
$E_0+h(t)dt\wedge\eta_i$, where $h$ is a smooth bump of integral one
vanishing near the chosen seam. The notation with $h=1$ below is its
piecewise polynomial path-model shorthand. Its cube is zero, as $E_0^2=0$ and the mixed
terms can be supported in disjoint strips. Independently, if $\alpha_i$
are the extra $\delta$-rows of the actual lifted monodromy, then
\[
 \alpha_1=-2G+U+K,\quad \alpha_2=3G+W-K,\quad\alpha_3=0,
 \qquad\eta_i=-\alpha_i+(A_i^T-I)(-U).
\]
Thus this representative is the actual principal-circle curvature class
up to the stated convention for its overall sign, not a class inferred
from its vertical part alone.

\paragraph{Geometric cap restrictions.}
On a finite cap the constant one-forms $G,K$ descend rationally to the
finite affine torus quotient. Untwisting the affine path by
$x\mapsto x-tv_i/m_i$ pulls a constant form $\zeta$ back to
$\zeta-dt\wedge\iota_{v_i/m_i}\zeta$. This is an actual pullback of forms,
so it preserves products. Consequently a decomposable representative
of $x$ on each finite cap restricts to
\[
 E_0+dt\wedge\beta_i,
 \qquad \beta_1=-\tfrac23K,\quad\beta_2=\tfrac12K.
\]
This convention uses the linear boundary gluing $(x,1)\sim(A_ix,0)$;
its affine cap map is $[x,t]\mapsto[x-tv_i/m_i]$. Thus its normal
meridian is $m_i\xi+v_i$. Reversing the angular convention changes
the final orientation sign, not the absolute cube. The sign also follows
from evaluating on the exact normal tori
$(-\omega(k,v_i),j_i(N_i k))$ in the companion integral calculation.
On each nodal cap, $G$ and $K$ extend as closed one-forms: in the marking
$a_1=w,a_2=w-\delta,b_1=u,b_2=g+u$, the quotient uses the transverse
primitive circle, and $G=b_2^*$: this is the transverse circle on the first cap and
the surviving longitude of the pinched factor on the second cap. The longitude form can be made constant
near the node by an exact, monodromy-invariant change. Its products with
$K$ give a cap representative restricting to $E_0$, so $\beta_3=\beta_4=0$.
All four cap representatives have square zero.

\paragraph{Boundary primitives.}
Put $b_i=\eta_i-\beta_i$.  In the piecewise polynomial notation, a one-form
$\Theta_i(t)=\lambda_i+t b_i$ is globally defined on the boundary mapping
torus whenever $(A_i^T-I)\lambda_i=b_i$; then
$d\Theta_i=dt\wedge b_i$ is precisely exterior minus cap curvature.
For smooth representatives replace $t\eta_i$ by $H(t)\eta_i$,
where $H'=h$ and $H(0)=0,H(1)=1$, retaining the term $-t\beta_i$.
Since $\beta_i\wedge E_0=0$, this change leaves every cube integral below
unchanged. The endpoint comparison is then smooth as well. One can take
\[
\begin{array}{c|c|c}
i& b_i&\lambda_i\\\hline
1&-4G-W+\tfrac23K&(-U+W)/3\\
2&-3G+U+\tfrac12K&(-U+W)/2\\
3&0&0
\end{array}
\]
On the last boundary the based inverse word gives three successive
transported segments
\[
 -C_2^TC_1^TA_2^T\eta_1=-2G+U+K,\quad
 -C_2^TC_1^T\eta_2=3G-U-K,\quad
 -C_2^T\eta_3=0.
\]
Their sum is $G=(C_2^T-I)W$, so start with $\lambda_4=W$ and add the
successive segment integrals. The only $W$ coefficient of $\Theta_4$
remains $1$, while the two nonzero $U$ increments are $+1,-1$.
The last-boundary contribution below therefore cancels exactly;
it is unchanged if the two nonzero segments are reversed.

\paragraph{Multiplicative collar calculation.}
Across a collar with coordinate $s$ from cap to exterior use the closed
form
\[
 \Omega_i=E_{\rm cap}+d(s\Theta_i)
          =E_{\rm cap}+s\,d\Theta_i+ds\wedge\Theta_i.
\]
At its endpoints it equals the chosen cap and exterior representatives.
Both $E_{\rm cap}^2$ and $(d\Theta_i)^2$ vanish. Hence
\[
 \Omega_i^3=6s\,ds\wedge\Theta_i\wedge E_0\wedge dt\wedge b_i.
\]
Integrating $s$ gives the factor $3$. With the algebraic boundary
orientation which records the coefficient of
$\Theta_i\wedge E_0\wedge b_i$ in $G\wedge U\wedge W\wedge K$,
the companion code uses fibre orientation $G\wedge U\wedge W\wedge K$
and obtains
\[
 3[\lambda_1E_0b_1]=2,\qquad
 3[\lambda_2E_0b_2]=-3,\qquad
 3[\Theta_3E_0b_3]=3[\Theta_4E_0b_4]=0.
\]
Thus the sum has absolute value $1$. For completeness, the common
overall orientation sign can be left unassigned: moving $dt$ past $\Theta_iE_0$ contributes a minus sign if
the collar orientation is $ds\wedge dt\wedge G\wedge U\wedge W\wedge K$.
The raw algebraic sum in the code is $-1$.

\paragraph{Explicit change to the inverse mapping-torus convention.}
The calculation above uses $(x,1)\sim(Ax,0)$. The convention
$(x,1)\sim(A^{-1}x,0)$ is obtained by reversing the peripheral
parameter and changing fibre coordinates by $A$ so that the marked fibre
at the basepoint is unchanged. Explicitly,
$[x,t]\mapsto[Ax,1-t]$ is the comparison and is the identity on the marked
$t=0$ fibre. The inverse-convention cap map
$[y,\theta]\mapsto[y+\theta v/m]$ becomes
$[Ax+(1-t)v/m]=[x-tv/m]$ after applying the inverse affine deck
transformation. Thus it is exactly the negative cap term used above.
The ribbon order and all fourth-boundary segments are reversed together.
The data transform as
\[
 A'=A^{-1},\quad \eta'=-A^{-T}\eta,\quad
 \beta'=-A^{-T}\beta=-\beta,\quad\lambda'=\lambda,
 \qquad b'=-A^{-T}b=(A^{-T}-I)\lambda.
\]
These are simultaneous changes of the actual forms, not merely changes
of the displayed matrices. Since $E_0$ is invariant and $\det A=1$,
\[
 [\lambda E_0b']=-[(A^T\lambda)E_0b]
                =-[(\lambda+b)E_0b]=-[\lambda E_0b].
\]
For the multi-segment last boundary use the full reversed primitive
$\Theta'(t)=A^{-T}\Theta(1-t)$; its initial value is again $\lambda$,
and the whole collar integral changes sign. Thus reversing the common
boundary convention multiplies all four contributions by the same sign.
The absolute value $1$ is independent of this choice.

Finally, $H^2(Q;\mathbb Q)$ is one-dimensional and restriction sends its
primitive generator to $2GK$, so the global closed form constructed
above represents the actual Euler class up to its independently fixed
sign. This proves the asserted absolute cube. The independently oriented
value $\ell=-4$ and the spin congruence $\ell\equiv4k\pmod{24}$ select
$k=-1$ from the two possibilities $\{1,-1\}$. The conclusion about
$\Theta_7$ is now B.6.19.\hfill$\square$
